\subsection{Elliptic surface volumes and order spectra}\label{sec:nef-surface-orders}

A coefficient that vanishes along one elliptic factor determines the harmonic mean in the volume formula. Finite-order zeros give fractional or logarithmic decay; flat coefficients allow the logarithmic slope to oscillate.

Let $E=\mathbb C/(\mathbb Z+\mathrm{i}\mathbb Z)$, write $X=E_z\times E_w$, and fix
\[
 \omega_1=\frac{\mathrm{i}}{2}\mathrm{d} z\wedge\mathrm{d}\bar z,\qquad
 \omega_2=\frac{\mathrm{i}}{2}\mathrm{d} w\wedge\mathrm{d}\bar w,\qquad
 \mu=\omega_1\wedge\omega_2,\qquad h=\omega_1+\omega_2.
\]
Both area forms have integral one. Let $a,q\geq0$ and $r$ be smooth functions of $z$ alone such that $|r|^2\leq aq$, and set
\[
 \theta=\frac{\mathrm{i}}{2}\bigl(a\,\mathrm{d} z\wedge\mathrm{d}\bar z+
 r\,\mathrm{d} z\wedge\mathrm{d}\bar w+\bar r\,\mathrm{d} w\wedge\mathrm{d}\bar z+
 q\,\mathrm{d} w\wedge\mathrm{d}\bar w\bigr).
\]
This is a smooth semipositive, hence nef, form.

\begin{proposition}\label{prop:nef-elliptic-exact-volume}
For $\epsilon>0$, put $g=q+\epsilon$ and $I=\int_Eg^{-1}\omega_1$. Then
\begin{equation}\label{eq:nef-elliptic-volume}
 \mathcal V_\mu(\theta+\epsilon h)
 =\frac{2}{I}\left(\int_E(a+\epsilon)\omega_1-
 \int_E\frac{|r|^2}{g}\omega_1\right).
\end{equation}
The infimum defining this volume is attained by a smooth positive Gauduchon metric. In particular, if $a=b>0$ and $r=0$, then $\mathcal V_\mu(\theta+\epsilon h)=2(b+\epsilon)/I$. If $|r|^2=aq$, then
\begin{equation}\label{eq:nef-rank-one-elliptic-volume}
 \mathcal V_\mu(\theta+\epsilon h)=\frac{2\epsilon(1+K)}{I},
 \qquad K=\int_E\frac{a}{q+\epsilon}\omega_1.
\end{equation}
\end{proposition}

\begin{proof}
Define
\[
 c=I^{-1}\left(\int_E(a+\epsilon)\omega_1-
 \int_E\frac{|r|^2}{g}\omega_1\right),\qquad
 A=\frac{c+|r|^2}{g}.
\]
The inequality $(a+\epsilon)g-|r|^2>0$ gives $c>0$, while the definition of $c$ makes $A-a-\epsilon$ have zero $\omega_1$-mean. Division of its nonconstant Fourier coefficients by the nonzero flat Laplace eigenvalues produces a smooth real function $u$ on $E$ with $\ddc u=(A-a-\epsilon)\omega_1$. Pull it back to $X$. The form $\Omega=\theta+\epsilon h+\ddc u$ has coefficient matrix $\left(\begin{smallmatrix}A&r\\\bar r&g\end{smallmatrix}\right)$, determinant $c$, and $\Omega^2=2c\mu$.

The metric $\gamma_*=\Omega/g$ has constant fiber coefficient $1$; its remaining coefficients depend only on $z$. Applying $\ddc$ to its base or mixed terms repeats $\mathrm{d} z$ or $\mathrm{d}\bar z$, so $\ddc\gamma_*=0$. It is smooth, positive, and pointwise conformal to $\Omega$. The arithmetic--geometric mean argument in \cref{thm:gaud-volume-dual-formula} and equality for $\gamma_*$ prove \eqref{eq:nef-elliptic-volume}. The two specializations follow by direct substitution.
\end{proof}

\begin{proposition}\label{prop:nef-diagonal-fractional-and-zero}\index{Nef degeneration!fractional and zero orders}
There are smooth semipositive nonclosed forms on the fixed surface $X$ with common order $1-1/m$ for each integer $m>1$. There is also one with
\[
 \mathcal V_\mu(\theta+\epsilon h)\sim\frac{2b}{\pi\log(1/\epsilon)}
 \quad(\epsilon\downarrow0),
\]
for any chosen $b>0$; it has boundary volume zero and both orders zero.
\end{proposition}

\begin{proof}
In \cref{prop:nef-elliptic-exact-volume} take $a=b>0$ and $r=0$. Choose $q$ smooth, positive outside $z_0\in E$, and equal to $|\zeta|^{2m}$ in a small flat coordinate disk about $z_0$. A positive smooth extension exists away from the disk. For $m=1$, polar coordinates give
\[
 I=\int_E\frac{\omega_1}{q+\epsilon}
   =\pi\log(1/\epsilon)+O(1).
\]
Equation \eqref{eq:nef-elliptic-volume} gives the displayed asymptotic. If $m>1$, the change $s=|\zeta|^2=\epsilon^{1/m}t$ gives
\[
 I\sim C_m\epsilon^{1/m-1},\qquad
 C_m=\pi\int_0^\infty\frac{\mathrm{d} t}{1+t^m}>0,
\]
and the common order is $1-1/m$. In the logarithmic case $\lambda_\mu(\theta)\geq0$ by semipositivity, while $\lambda_\mu(\theta)\leq\lambda_\mu(\theta+\epsilon h)\to0$ by monotonicity. Hence $\mathcal V_\mu(\theta)=0$. The quotient of the displayed volume by $\epsilon^\delta$ diverges for every $\delta>0$. Since $\mathrm{d}\theta=\mathrm{d} q\wedge\omega_2\ne0$, these backgrounds are nonclosed.
\end{proof}

\begin{proposition}\label{prop:nef-oscillatory-orders}\index{Nef degeneration!unequal orders}
For every $0<\alpha<\beta<1$ there is a smooth semipositive nonclosed background on the same $(X,\mu,h)$ with
\[
 \bigl(\underline a(\theta),\overline a(\theta)\bigr)
 =(1+\alpha,1+\beta).
\]
For every $0<\sigma<1$ there is such a background with common order $1+\sigma$.
\end{proposition}

\begin{proof}
For all sufficiently large $t$, define
\[
 p(t)=\frac{\alpha+\beta}{2}
 +\frac{\beta-\alpha}{2}\sin(\log\log t),\qquad P(t)=tp(t).
\]
Choose $T$ large enough that $0<c_0\leq P'(t)\leq d_0<1$ for $t\geq T$; this follows from
\[
 P'(t)=p(t)+\frac{\beta-\alpha}{2\log t}\cos(\log\log t).
\]
On $E\setminus\{z_0\}$ choose a smooth $t(z)\geq T$ equal to $|\zeta|^{-2}$ near $z_0$ and bounded outside every neighborhood of $z_0$. Put
\[
 q=\mathrm{e}^{-t},\qquad a=\mathrm{e}^{-P(t)},\qquad
 r=\mathrm{e}^{-(P(t)+t)/2},
\]
and set all three equal to zero at $z_0$. These functions are smooth and flat there: $P(t)\geq\alpha t$ and every fixed derivative of $P$ grows at most polynomially in $t$, so the chain rule bounds each coordinate derivative by a polynomial in $|\zeta|^{-1}$ times $\mathrm{e}^{-c/|\zeta|^2}$. We have $|r|^2=aq$, so the resulting $\theta$ is semipositive. It is nonclosed: the component $\mathrm{d} q\wedge\omega_2$ of $\mathrm{d}\theta$ has two fiber differentials, which derivatives of the remaining terms cannot cancel.

Write $\epsilon=\mathrm{e}^{-L}$. Choose a fixed integration endpoint $T_0\geq T$ inside the coordinate disk, with the transition to the bounded global $t(z)$ outside a smaller disk. The complementary contributions to $I$ and $K$ in \eqref{eq:nef-rank-one-elliptic-volume} are $O(1)$, while polar coordinates on the disk give
\[
 I_{\mathrm{loc}}=\pi\int_{T_0}^\infty
 \frac{t^{-2}\,\mathrm{d} t}{\mathrm{e}^{-t}+\mathrm{e}^{-L}},\qquad
 K_{\mathrm{loc}}=\pi\int_{T_0}^\infty
 \frac{\mathrm{e}^{-P(t)}t^{-2}\,\mathrm{d} t}{\mathrm{e}^{-t}+\mathrm{e}^{-L}}.
\]
Integrating $I$ over $t\geq L$ gives a lower bound by a fixed multiple of $\mathrm{e}^L/L$. For the upper bound, use $(\mathrm{e}^{-t}+\mathrm{e}^{-L})^{-1}\leq\mathrm{e}^t$ for $t\leq L$ and at most $\mathrm{e}^L$ for $t\geq L$, then split $[T_0,L]$ at $L/2$. Thus
\begin{equation}\label{eq:nef-elliptic-I-asymptotic}
 \log I=L-\log L+O(1).
\end{equation}
The derivatives of $P$ and $t-P(t)$ are positive and bounded away from zero. Splitting the $K$-integral at $L$ therefore gives $K\leq C\mathrm{e}^{L-P(L)}$. On $[L,L+1]$ its denominator is at most $2\mathrm{e}^{-L}$ and $P(t)\leq P(L)+1$, so
\[
 c_1L^{-2}\mathrm{e}^{L-P(L)}\leq K\leq c_2\mathrm{e}^{L-P(L)}.
\]
The complementary $O(1)$ term is absorbed because $P(L)\leq\beta L$. Consequently $\log K=L-P(L)+O(\log L)$ and $K\to\infty$. Combining this with \eqref{eq:nef-rank-one-elliptic-volume} and \eqref{eq:nef-elliptic-I-asymptotic} yields
\begin{equation}\label{eq:nef-oscillatory-log-volume}
 \log \mathcal V_\mu(\theta+\mathrm{e}^{-L}h)=-L-P(L)+O(\log L),
 \qquad
 \frac{\log \mathcal V_\mu(\theta+\epsilon h)}{\log\epsilon}
 =1+p(L)+o(1).
\end{equation}
The sine in $p(L)$ has limits $-1$ and $1$ along sequences, proving the first assertion. If instead $P(t)=\sigma t$, the same smoothness and integral estimates give common order $1+\sigma$. In particular $\sigma=\sqrt2-1$ gives an irrational common order.
\end{proof}
